% key feedback points
% active voice
% straightforward, like how it would normally be said out loud
\chapter{Discussion}\label{ch:discussion}
\id{Use active voice everywhere: In this chapter, I discuss the key takeaways of my Master's research and identify future directions for prospective researchers.} This chapter reflects on the reliability model introduced in \chref{ch:reliable-sots} and the demonstrative evaluation presented in \chref{ch:demonstration}. It first analyzes \id{no analysis at this point. analysis was in earlier chapters. here, you elaborate on stuff} the utility of the approach (\secref{sec:discussion-utility}), then discusses implications for SoTS engineering practice (\secref{sec:discussion-implications}), and, finally, outlines directions for prospective researchers (\secref{sec:discussion-future}).

\section{Utility of the Approach}\label{sec:discussion-utility}
TODO - ONCE EVALUATION IS DONE
\subsection{Model-Driven Coordination and Enforcement}
TODO - ONCE EVALUATION IS DONE



\subsection{Graceful Degradation and Adaptive Reliability}
TODO - ONCE EVALUATION IS DONE
 


\subsection{Model-to-Code Integration and Complexity Reduction}
TODO - ONCE EVALUATION IS DONE



\section{Implications for SoTS Engineering Practice}\label{sec:discussion-implications}
The proposed approach has several implications for how SoTS are to be engineered and managed, particularly in structuring system-level coordination across independently developed constituents.


\subsection{Statecharts for Engineering Process Maturity}

\id{So, this title doesn't sound like an implication. When I read "implications," I expect clear statements, either in the title or in the very first sentence. Is the implication that ``statecharts for engineering process maturity''? This sounds more like a technique rather than an implication. Rephrase these points please.}

Lifecycle-based control advances engineering maturity by making system-level interaction explicit, structured, and enforceable. At early stages, lifecycle models provide a shared representation of participation, improving consistency across teams. At higher maturity levels, they support formalised coordination policies, integration into toolchains for automated code generation and validation, and runtime enforcement of behaviour. This aligns with frameworks such as Capability Maturity Model Integration (CMMI) \cite{paulk1991capability}, where maturity progresses from informal practices to defined and quantitatively controlled processes.

Lifecycle models can be incorporated into existing engineering workflows as managed artefacts, similar to requirements or architectural models, enabling versioning, review, and validation. At higher maturity levels, they also support runtime monitoring of system behaviour, allowing assessment of compliance with coordination policies and iterative refinement. Adopting this approach requires capabilities in modelling and system-level design, along with governance to manage evolving lifecycle definitions. It may also introduce roles such as SoTS architects or model engineers responsible for maintaining these models. As shared reference points, executable lifecycle models improve communication, reduce ambiguity, and support more controlled engineering processes.


\subsection{Statecharts as Enforceable Contracts in SoTS}
Lifecycle statecharts can be interpreted as executable coordination contracts between DT constituents in the SoTS. States define the conditions for valid participation, while transitions capture permitted behavioural changes, making assumptions and guarantees explicit. As these contracts are enforced through statechart execution, they regulate participation, contributions, and behaviour at the SoTS level rather than within individual constituents. This allows each DT to rely on consistent participation rules, assuming all others are controlled by the same lifecycle constraints. This aligns with contract-based design \cite{benveniste2012contracts}, where behaviour is defined through assumptions and guarantees, and is realized here through statecharts that both describe and enforce valid interaction across independently evolving constituents. Previous work shows that contract specifications can be transformed into state-based models \cite{bae2016systematic}, enabling runtime enforcement through states and transitions.

This perspective also supports reasoning about safety and system evolution. Disturbances and degraded conditions represent situations where guarantees can no longer be maintained, triggering constrained transitions such as restriction or exit to preserve system-level behaviour \cite{bozzano2014formal, sangiovanni2012taming}. At the same time, lifecycle contracts establish a stable interaction boundary, allowing systems to evolve independently while maintaining  compatibility at the SoTS level \cite{waschle2021contract}.


\subsection{Separation of Concerns and Privacy in SoTS}
The proposed approach separates constituent internals from system-level coordination by externalising coordination logic into lifecycle models. DT constituents participate through well-defined states and transitions without exposing internal behaviour or implementation details, while interaction is defined through shared lifecycle constraints rather than embedded within individual components. This keeps coordination at the SoTS level and allows constituents to focus on local functionality. As a result, constituents can evolve independently---being modified, replaced, or reconfigured---without global redesign, provided they conform to shared lifecycle definitions. This aligns with operational and managerial independence in SoS \cite{maier1998architecting} and architectural principles that limit exposure to stakeholder-relevant concerns \cite{iso42010}, enabling scalable integration of heterogeneous DTs while maintaining system-level consistency.

This separation also limits information exposure and reduces integration and cognitive complexity. Interaction is mediated through lifecycle states, ensuring only participation-relevant information is exchanged, while internal models and logic remain encapsulated within each DT. This aligns with privacy risk management perspectives such as the NIST Privacy Framework \cite{enterprise2020nist}, which emphasise addressing privacy risks through system design across the data lifecycle. By constraining participation, lifecycle-based coordination reduces unnecessary data flows and unintended exposure---particularly in SoS contexts, where interactions between independent systems can introduce unexpected security, privacy, and trust risks \cite{olivero2024systematic}. It also shifts engineering effort: SoTS-level developers focus on disturbance management and system guarantees, while service-level developers focus on local behaviour without reasoning about broader system dynamics. Lifecycle-based coordination thus provides a unified mechanism for managing interaction, evolution, and information exposure in SoTS environments.


\section{Future Directions}\label{sec:discussion-future}
The proposed approach establishes a foundation for state-based \id{I don't know what you mean and I don't know why are you underselling the contributions by this overly specific adjective.} coordination in SoTS. However, several directions remain for extending the model's expressiveness and adaptability.

\subsection{Extending the Orthogonal Dimensions in the Model}
\id{Orthogonal regions, not dimensions. But more importantly: you want to talk about additional considerations that, in turn, can be captured as orthogonal dimensions. You need to turn this around and talk about those additional concerns and then say that btw, these can be orthogonal regions}
The lifecycle-based model can be extended by leveraging the ability of statecharts to encode multiple orthogonal regions. Each region represents a distinct concern influencing participation. While the current model captures belonging and health, additional dimensions such as trust, operational readiness, data integrity, security, and quality-of-service (QoS) attributes \cite{o2007quality} can be incorporated as independent but interacting regions. These concerns already play a central role in digital twin environments when determining which systems to interact with. They can also be encoded at the SoTS level. Trust, for example, evolves continuously. It can be evaluated based on system interactions, historical behaviour, and recommendations from other constituents \cite{li2023trust}. It directly influences participation by enabling or restricting access to system-level interactions.

However, introducing multiple dimensions raises new challenges. Constraints across regions may conflict, requiring prioritization or coordination policies to resolve trade-offs between assurance, performance, and system capability. In evolving SoTS, constituent systems change independently. As a result, relationships must be preserved or adapted over time \cite{carney2005topics}, making these conflicts more pronounced.  In cross-organizational settings, this is further complicated by the need for shared semantics. Digital twins must agree not only on interaction protocols, but also on how dimensions such as trust, readiness, and data validity are defined and evaluated. Embedding these dimensions into the state representation transforms lifecycle models into multi-dimensional coordination structures.

\subsection{Adaptive Fault Tolerance in SoTS}
In the presented model, the highest level of reliability introduces a compensation component that directly implements a fault tolerance strategy under degraded conditions \cite{hanmer2013patterns}. This can be extended to support dynamic selection based on system context, allowing behaviour to be refined at runtime while remaining within lifecycle constraints. In adaptive systems, fault tolerance must account for evolving conditions, motivating context-sensitive strategies \cite{de2022systematic}. Fault tolerance techniques inherently involve trade-offs with performance, resource usage, and timeliness. For instance, replication and recovery improve reliability but may introduce latency and unpredictability, conflicting with real-time requirements \cite{narasimhan2002trade}. Similarly, resource-level adaptations such as frequency scaling and task reallocation expose trade-offs between reliability and efficiency \cite{panerati2015trading}. Dynamic mechanisms enable systems to navigate these trade-offs by selecting strategies aligned with current conditions, improving both efficiency and resilience compared to static approaches \cite{ballesteros2022infrastructure}.

This can be realized in the current model by introducing an orthogonal decision region within the statechart to manage strategy selection, operating independently of the belonging and health dimensions while respecting lifecycle constraints. Strategy selection may be guided by policies or utility-based mechanisms \cite{avizienis2004basic}, and further refined through learning-based approaches such as reinforcement learning \cite{dong2023deep, zhang2020fault}. This extends the lifecycle model from enforcing allowable configurations to enabling adaptive decision-making within those constraints, improving flexibility while preserving reliability guarantees.


\subsection{Verification and Analysis of Statechart-Driven Behaviour}
\id{``Verification of SoTS.'' Cut back this verbose language. It really hurts this chapter. Just write like you would normally speak. If there's anything related to your work, it must be verification of SoTS, not statechart-driven behavior. A: this expression doesn't make sense. B: behaviour of what? My car? My phone? No, behaviour of SoTS. Explain that with explicitly modeled behaviour, this exercise is now easy-peasy. That will drive this discussion home.}
A key direction for future work is to establish formal guarantees over statechart-driven behaviour in SoTS, where reliability depends on constraining disturbance-causing configurations. The proposed lifecycle statechart model encodes permitted configurations and transitions during operation (\chref{ch:reliable-sots}), reducing the likelihood of disturbance propagation. This enables statecharts to be treated as analyzable specifications of system behaviour. This supports formal reasoning over the state space and system evolution. Properties such as safety (forbidden configurations are never reached), liveness (desired progression is ensured), and invariants over belonging and health can be analysed. It also enables verification of reachability, transition correctness, and following of reliability constraints, ensuring behaviour remains within defined bounds \cite{latella1999towards, sheldon2002validation}.

To support deeper analysis, model-driven engineering (MDE) techniques can be used to derive additional representations of statechart-driven behaviour. In MDE, transformations systematically generate alternative models from a primary representation \cite{biehl2010literature}. This allows behaviour to be defined once while supporting multiple analytical views \cite{bezivin2006model}. While statecharts are well suited for modelling and enforcing runtime behaviour, they do not explicitly capture properties such as resource usage, synchronization structure, or global system conditions. These aspects are difficult to analyse directly within the formalism.m Through model-to-model transformations, statecharts can be converted into representations such as Petri nets \cite{cortadella1995synthesizing, hei2011automatic}, which provide explicit modelling of concurrency, synchronization, and resource flow. This enables analysis of properties that are not directly visible in the original model, including deadlocks, unreachable states, and resource contention. The resulting insights can inform refinements to the statechart, such as introducing additional constraints or guarded transitions to enforce safe behaviour. This approach supports systematic reuse of behavioural models while maintaining consistency across representations. It strengthens the integration between design, execution, and verification, enabling more reliable statechart-driven behaviour in SoTS.

